\section{Results and Conditional Interpretation}
\subsection{First-Contact Loss Depends on the Declared Horizon}
Both laws completed 48/48 healthy RF-0 tasks; Fig.~\ref{fig:v10f-transfer}(a) shows paired losses. Fig.~\ref{fig:v10f-transfer}(d) rescores the same 50\%-deletion trajectories at 900/1000\,ms across all three RFs. RF-0 has nine additional P contacts in the final 100\,ms (latest 999.962\,ms); law-discordant RF-A/B contacts also occur after 900\,ms. These descriptive rescorings add neither samples nor confirmatory tests; RF-0's repeated endpoints were registered.

The prespecified secondary 1-s R-minus-P net-loss contrast (Fig.~\ref{fig:v10f-transfer}(b)) was $-22.9$\,pp (95\% CI $[-33.3,-12.5]$; Holm $p=0.00195$). RF-A/B each yielded $-4.2$\,pp, CI $[-12.5,0]$; all four new map/dose Holm $p$ values are 1. Each new 50\% contrast has only one nonzero cluster. These estimates establish neither equivalence nor an isolated RF effect; no between-map test is claimed. The baseline 25\% contrast remains uncertain ($p=0.50$). No 25\% loss occurs in A/B; twelve event-free clusters still allow a one-sided 95\% upper bound of 22.1\% for any loss in a new identically defined two-bearing cluster.

All 576 records triggered. Post-hoc healthy--50\%-deletion delays averaged 80.0/57.5/57.5\,ms in RF-0/A/B (SD 42.9/46.6/29.7\,ms across 48/24/24 bearing pairs; 24/12/12 independent seeds, P/R duplicates counted once). Trigger means first post-cue crossing by the eight-cell five-frame mean at 20-ms resolution. All 288 P/R pairs shared trigger/burst schedules under pre-trigger gating, yet 15 had different contact outcomes. Both RF-0 25\% P contacts occurred without delay: timing alone neither determines outcome across laws nor accounts for every contact.

Among 543 contact-free burst completions with follow-up, all 3,114 decisions at/after burst end and before contact/$H$ had zero eight-cell features and gated turn; 474 records retained whole-network spikes. Their 13 later contacts occurred during this observed walking regime. This post-hoc stage result contextualizes burst-end geometry and finite follow-up; it is not a universal silence rule or an isolated delay mechanism. Four prior contacts, 17 ongoing bursts and 12 zero-follow-up records remain outside this risk set.

In Fig.~\ref{fig:v10f-transfer}(c), paired healthy minimum-clearance points and vertical median marks show unequal reserve: RF-0 P/R medians were 33.1/41.4\,mm (400-ms absolute heading $88.3^\circ/100^\circ$). Reserve, burst-end state and late follow-up are possible contributors, not isolated mechanisms; RF calibration and cohorts also vary together.

\subsection{Frozen Scalar Warnings Miss Subsequent Contacts}
\WarningFigure
Across healthy/25\%/50\% conditions, the command warning detected 2/13 RF-0 contacts, 2/3 in A and 0/1 in B (\warningRef(a)); successful-task false alarms were 51/275, 42/141 and 24/143 (\warningRef(b)). Exact-count re-scoring preserves all historical labels and CIs; rate-first rounding would falsely flag 19 RF-0 successes tied at 14 counts. The upper-only activity rule detected none. Exact-state forecasts detected all observed contacts but falsely flagged 57/275, 25/141 and 33/143 successes; baseline schedule-only false alarms were 84/275. Their geometry/actuation information precludes a neural-only comparison.

RF-0/A/B sensitivities have only 12/2/1 failure clusters. Of 20,000 resamples, 2,217/7,099 lack a positive class in A/B; those draws are omitted only for the corresponding conditional rate. Zero-width empirical intervals do not establish zero population risk. Own-stop/common-duration medians coincide in these campaigns (maximum paired difference below 0.006), so the historical large exposure artifact does not recur. P/50\% median 400-ms/1-s activity ratios remain 0.981/1.505, 0.961/1.212 and 0.956/1.139. A separate eight-cluster $\pm45^\circ$ test flags 22/32 healthy activity records and retains an uncertain law contrast. These results bound the tested frozen criteria, not all neural monitoring.
